\chapter{FD2NN Fundamental-Limit Discussion}

\section{What This Chapter Does and Does Not Claim}

This chapter is deliberately conservative. It does not present a strict
impossibility theorem for every phase-only FD2NN. Instead, it combines three
observations:
\begin{enumerate}
  \item turbulence PSDs describe second-order statistics, not realization-level
  phase;
  \item a static Fourier-plane mask cannot implement a realization-specific
  inverse phase;
  \item the corrected \texttt{0405} 4f sweep produces only a near-tie at its best
  tested geometry.
\end{enumerate}

\section{PSD versus Realization}

Let a Fourier-domain field be written as
\begin{equation}
\hat{\psi}(\kappa)=|\hat{\psi}(\kappa)|\,e^{j\theta(\kappa)}.
\end{equation}
The PSD determines $|\hat{\psi}(\kappa)|^2$ statistically, but it does not
specify the realization-dependent phase $\theta(\kappa)$. A static Fourier mask
has the form
\begin{equation}
H(\kappa)=e^{j\Delta\phi(\kappa)},
\end{equation}
whereas an ideal realization-specific correction would require
\begin{equation}
H_{\mathrm{ideal}}(\kappa)=e^{-j\theta_{\mathrm{turb}}(\kappa)}.
\end{equation}
This is the structural mismatch: a fixed mask cannot equal a realization-varying
inverse phase law.

The same point can be seen in the simplest one-mode toy example. Suppose a
single Fourier coefficient carries a turbulent phase $\theta$ and the static
compensator applies a fixed phase shift $\Delta \phi$. The corrected factor is
\begin{equation}
e^{j(\theta + \Delta \phi)}.
\end{equation}
If the turbulent phase is broadly distributed, then
\begin{equation}
\mathbb{E}[e^{j(\theta + \Delta \phi)}]
=
e^{j\Delta\phi}\,\mathbb{E}[e^{j\theta}]
\approx 0,
\end{equation}
so no single fixed phase shift can sharpen the expectation. The missing
information is the realization-specific phase, not the choice of a better fixed
mask.

\section{Toy Wiener Intuition}

The project presentation material includes a single-scalar Wiener-style argument
for why a fixed Fourier correction tends to collapse toward a trivial average
operator. That argument is useful as intuition, but it is not a direct proof
for a multilayer phase-only FD2NN. It should therefore be framed as an
interpretive aid rather than a theorem.

\section{Empirical Evidence from the 4f Sweep}

The corrected \texttt{0405} sweep is consistent with the structural mismatch above. In
the tested geometry, the best 4f setting raises the focal PIB by only
$+0.0104$, while spatial D2NN branches achieve much larger detector-side gains.
The safe statement is not ``FD2NN can never work.'' It is narrower: in the
tested 4f geometry, static Fourier-plane phase correction provides at most a
very weak detector-side gain.

\section{The $f$--$z$ Trade-off}

The 4f architecture also inherits a geometry dilemma. Decreasing $f$ increases
mixing strength but worsens sampling pressure and defocus sensitivity.
Increasing $f$ improves physical credibility but weakens the very coupling that
would be needed to build a useful correction. The tested sweep never escaped
this trade-off.

\begin{discovery}
The 4f FD2NN results do not justify a universal impossibility claim. They do
justify a strong practical warning: static Fourier-plane phase masks appear far
less effective for random turbulence cleanup than the best spatial
detector-centric D2NN designs studied in the same campaign.
\end{discovery}

\section{Transition to the Final Conclusion}

The report therefore closes with a nuanced picture. Static diffractive optics
can matter at the detector, but the useful gains emerge from detector-centric
spatial shaping, not from field-level restoration or strong Fourier-plane
correction.
